\chapter{Relative Spec}
\label{chap:Picard_RelativeSpec}
% archon:covers AlgebraicJacobian/Picard/RelativeSpec.lean

% STRATEGY NOTE (iter-172). This chapter is the foundational A.1.a sub-row of
% Route A (positive-genus arm of nonempty_jacobianWitness). It packages the
% relative-spectrum functor \underline{\Spec}_X(\mathcal{A}) : \mathrm{QcohAlg}(X)^{op}
% \to \mathrm{Sch}/X used by the relative Picard functor on a product C \times_k T.
% Mathlib has the absolute Spec functor and affine gluing (Scheme.GlueData,
% AffineScheme.glueOpens) but no general RelativeSpec construction; this chapter is
% the prover-ready specification for the new file AlgebraicJacobian/Picard/RelativeSpec.lean
% which iter-173 (file-skeleton lane) will scaffold.

\section{Setup and motivation}
\label{sec:relspec_setup}

The Route~A construction of \(J = \Pic^0_{C/k}\) as a \(g\)-dimensional abelian variety is
organised around the relative Picard functor
\(\Pic^\sharp_{C/k}(T) = \Pic(C \times_k T) / \pi_T^* \Pic(T)\)
on the category of \(k\)-schemes. Representing \(\Pic^\sharp_{C/k}\) requires the line-bundle
pullback construction on a relative curve \(C \times_k T \to T\), which in turn rests on a
\emph{relative scheme} construction on a quasi-coherent sheaf of \(\mathcal{O}_X\)-algebras
\(\mathcal{A}\), the \emph{relative spectrum}
\(\pi : \underline{\Spec}_X(\mathcal{A}) \to X\). This chapter packages the construction
of \(\underline{\Spec}_X(-)\) and its core properties --- universal property and affine
base --- as the foundational A.1.a building block. The relative Picard
functor itself (A.1.c), the line-bundle pullback functor on a product (A.1.b), and the
identity-component / dimension theory of \(J\) (A.3) live in separate sibling chapters
(\S~``Out of scope'' below).

% SOURCE: [Stacks Project], Section ``Relative spectrum via glueing'' and ``Relative
% spectrum as a functor'' (tags 01LL, 01LO, 01LQ, 01LR, 01LS, available locally at
% references/stacks-constructions.tex L309-L720). The Hartshorne reference is
% [Hartshorne], II Exercise 5.17 (read from references/hartshorne-algebraic-geometry.pdf).

\section{Quasi-coherent sheaves of algebras}
\label{sec:qc_sheaf_of_algebras}

\begin{definition}\leanok
[Quasi-coherent \(\mathcal{O}_X\)-algebra]
  \label{def:qc_sheaf_of_algebras}
  \lean{AlgebraicGeometry.Scheme.QcohAlgebra}
  % SOURCE: [Stacks Project], Section 01LL ``Relative spectrum via glueing'',
  % Situation ``situation-relative-spec''
  % (read from references/stacks-constructions.tex, L312-L318).
  % SOURCE QUOTE:
  % "Here \(S\) is a scheme, and \(\mathcal{A}\) is a quasi-coherent
  %  \(\mathcal{O}_S\)-algebra. This means that \(\mathcal{A}\) is a
  %  sheaf of \(\mathcal{O}_S\)-algebras which is quasi-coherent as an
  %  \(\mathcal{O}_S\)-module."
  \textit{Source: [Stacks Project], tag 01LL (situation-relative-spec).}
  Let \(X\) be a scheme. A \emph{quasi-coherent sheaf of \(\mathcal{O}_X\)-algebras} is a
  sheaf \(\mathcal{A}\) of \(\mathcal{O}_X\)-algebras (i.e. an \(\mathcal{O}_X\)-module
  equipped with \(\mathcal{O}_X\)-linear multiplication and unit morphisms satisfying
  associativity, commutativity, and the unit law as sheaf morphisms) whose underlying
  \(\mathcal{O}_X\)-module is quasi-coherent. We write \(\mathrm{QcohAlg}(X)\) for the
  category of quasi-coherent sheaves of \(\mathcal{O}_X\)-algebras and
  \(\mathcal{O}_X\)-algebra morphisms.
\end{definition}

% NOTE (iter-179 review): supersedes the iter-174 carrier note. The iter-179
% Block A refactor (analogies/relative-spec-encoding.md) upgrades the carrier
% to a 3-field structure
%     `structure QcohAlgebra (X : Scheme.{u}) where`
%     `  sheaf        : TopCat.Sheaf CommRingCat.{u} X.toPresheafedSpace`
%     `  unit         : X.sheaf ⟶ sheaf`
%     `  coequifibered : NatTrans.Coequifibered`
%     `    (Functor.whiskerLeft (AffineZariskiSite.toOpensFunctor X).op unit.hom)`
% The new third field IS the iter-175+ overlay the prior note flagged: Mathlib
% commit `b80f227`'s Stacks-01LL `NatTrans.Coequifibered` predicate is strictly
% weaker than the full `SheafOfModules.IsQuasicoherent` claimed by Definition
% 1.1's prose but is equivalent for the geometric purpose (constructing the
% relative-spectrum scheme) via the dense-subsite equivalence
% `AffineZariskiSite.sheafEquiv`. This is the predicate Mathlib's
% `AffineZariskiSite.relativeGluingData` consumes, so the iter-179
% `RelativeSpec`/`structureMorphism` bodies (now real Mathlib values) close
% against it.

\section{The relative-spectrum scheme}
\label{sec:relspec_scheme}

The construction is given affine-locally: on every affine open \(U = \Spec R \subseteq X\) one
sets \(\pi^{-1}(U) := \Spec(\mathcal{A}(U))\) regarded as an \(R\)-algebra, and the local pieces
glue compatibly because \(\mathcal{A}\) is quasi-coherent (the transition isomorphism
\(\mathcal{A}(U) \otimes_R S \xrightarrow{\sim} \mathcal{A}(V)\) for an inclusion
\(V = \Spec S \subseteq U\) gives an open immersion of the corresponding Specs).

\begin{theorem}\leanok
[Relative spectrum exists]
  \label{thm:relative_spec_exists}
  \lean{AlgebraicGeometry.Scheme.RelativeSpec}
  \uses{def:qc_sheaf_of_algebras}

  % SOURCE: [Stacks Project], tag 01LQ (lemma-glue-relative-spec)
  % (read from references/stacks-constructions.tex, L381-L405).
  % SOURCE QUOTE:
  % "In Situation REF.
  %  There exists a morphism of schemes
  %  \(\)
  %  \pi : \underline{\Spec}_S(\mathcal{A}) \longrightarrow S
  %  \(\)
  %  with the following properties:
  %  \begin{enumerate}
  %  \item for every affine open \(U \subset S\) there exists an isomorphism
  %  \(i_U : \pi^{-1}(U) \to \Spec(\mathcal{A}(U))\) over \(U\), and
  %  \item for \(U \subset U' \subset S\) affine open the composition
  %  \(\)
  %  \xymatrix{
  %  \Spec(\mathcal{A}(U)) \ar[r]^{i_U^{-1}} &
  %  \pi^{-1}(U) \ar[rr]^{inclusion} & &
  %  \pi^{-1}(U') \ar[r]^{i_{U'}} &
  %  \Spec(\mathcal{A}(U'))
  %  }
  %  \(\)
  %  is the open immersion of Lemma REF above.
  %  \end{enumerate}
  %  Moreover, \(\underline{\Spec}_S(\mathcal{A})\)
  %  is unique up to unique isomorphism over \(S\)."
  \textit{Source: [Stacks Project], tag 01LQ (lemma-glue-relative-spec);
  cf.\ [Hartshorne], II~Ex.~5.17(a).}
  Let \(X\) be a scheme and \(\mathcal{A}\) a quasi-coherent \(\mathcal{O}_X\)-algebra
  (\cref{def:qc_sheaf_of_algebras}). There exists a scheme
  \(\underline{\Spec}_X(\mathcal{A})\) together with an affine morphism
  \(\pi : \underline{\Spec}_X(\mathcal{A}) \to X\) such that
  \begin{enumerate}
    \item for every affine open \(U \subseteq X\), there is an isomorphism
      \(i_U : \pi^{-1}(U) \xrightarrow{\sim} \Spec(\mathcal{A}(U))\) over \(U\);
    \item for \(U \subseteq U' \subseteq X\) both affine open, the composite
      $\Spec(\mathcal{A}(U)) \xrightarrow{i_U^{-1}} \pi^{-1}(U) \hookrightarrow
      \pi^{-1}(U') \xrightarrow{i_{U'}} \Spec(\mathcal{A}(U'))$
      coincides with the open immersion induced by the restriction
      \(\mathcal{A}(U') \to \mathcal{A}(U)\).
  \end{enumerate}
  The pair \((\underline{\Spec}_X(\mathcal{A}), \pi)\) is unique up to unique isomorphism
  over \(X\).
\end{theorem}

% SOURCE QUOTE PROOF: "Follows immediately from
% Lemmas REF,
% REF, and
% REF.
% Uniqueness is stated in the last sentence of
% Lemma REF."
\begin{proof}
  
  
\leanok
  By the open-cover gluing principle for schemes (Mathlib
  \texttt{Scheme.GlueData} / \texttt{AffineScheme.glueOpens}) it suffices to specify
  affine local pieces \(\Spec(\mathcal{A}(U))\) for each affine open \(U \subseteq X\) and
  open-immersion transition data
  \(\Spec(\mathcal{A}(U)) \hookrightarrow \Spec(\mathcal{A}(U'))\) for \(U \subseteq U'\)
  satisfying the cocycle condition. The transition morphisms come from the ring maps
  \(\mathcal{A}(U') \to \mathcal{A}(U)\); quasi-coherence yields the cartesian square
  \(\Spec(\mathcal{A}(U)) = U \times_{U'} \Spec(\mathcal{A}(U'))\), so each transition is an
  open immersion. Transitivity (Stacks lemma-transitive-spec) gives the cocycle condition.
  Uniqueness of the glued scheme is the standard property of the gluing functor.
\end{proof}

\begin{definition}

\leanok
[Structure morphism of the relative spectrum]
  \label{def:relspec_structure_morphism}
  \lean{AlgebraicGeometry.Scheme.RelativeSpec.structureMorphism}
  \uses{thm:relative_spec_exists, def:qc_sheaf_of_algebras}
  Let \(X\) be a scheme and \(\mathcal{A}\) a quasi-coherent \(\mathcal{O}_X\)-algebra.
  The \emph{structure morphism}
  \(\pi : \underline{\Spec}_X(\mathcal{A}) \to X\) is the canonical projection of the
  relative spectrum to its base, obtained as the colimit descent of the affine-local
  structure morphisms \(\Spec(\mathcal{A}(U)) \to U\) over the directed affine open
  cover of \(X\). It is the morphism whose affineness and base-change behaviour are the
  subject of the universal-property and base-change results below.
\end{definition}

\begin{proof}

\leanok
  The structure morphism is the base morphism of the relative gluing data of
  \(\mathcal{A}\); it is the projection underlying \cref{thm:relative_spec_exists}.
\end{proof}

\begin{theorem}\leanok
[Universal property of the relative spectrum]
  \label{thm:relative_spec_univ}
  \lean{AlgebraicGeometry.Scheme.RelativeSpec.UniversalProperty}
  % NOTE iter-283: removed the spurious `thm:relative_spec_base_change` from this
  % \uses. The universal property is the PRIMITIVE statement, proved by reducing
  % to the affine-base case and gluing (Lean body of `UniversalProperty` uses
  % only `isAffineHom_of_forall_exists_isAffineOpen`/`relativeGluingData`, never
  % `base_change`). Base change is a DOWNSTREAM corollary
  % (thm:relative_spec_base_change \uses thm:relative_spec_univ), so listing it
  % here created a univ <-> base_change dependency cycle that crashed
  % `leanblueprint web`.
  \uses{thm:relative_spec_exists, thm:relative_spec_affine_base,
    def:relspec_structure_morphism}

  % NOTE (iter-173 review): the iter-173 file-skeleton encodes this as the
  % structural consequence ``IsAffineHom (structureMorphism 𝒜)'' rather than
  % the natural-bijection statement pinned by this prose. Iter-174+ must
  % upgrade the Lean type to a ``CategoryTheory.Functor.RepresentableBy''
  % witness (or equivalent) before any consumer file relies on this theorem.
  % SOURCE: [Stacks Project], tag 01LQ continuation: Section ``Relative spectrum as
  % a functor'', functor F and Lemma lemma-spec
  % (read from references/stacks-constructions.tex, L427-L466 + L547-L551).
  % SOURCE QUOTE:
  % "For any \(f : T \to S\) the pullback
  %  \(f^*\mathcal{A}\) is a quasi-coherent sheaf of \(\mathcal{O}_T\)-algebras.
  %  We are going to consider pairs \((f : T \to S, \varphi)\) where
  %  \(f\) is a morphism of schemes and \(\varphi : f^*\mathcal{A} \to \mathcal{O}_T\)
  %  is a morphism of \(\mathcal{O}_T\)-algebras. Note that this is the
  %  same as giving a \(f^{-1}\mathcal{O}_S\)-algebra homomorphism
  %  \(\varphi : f^{-1}\mathcal{A} \to \mathcal{O}_T\), see
  %  Sheaves, Lemma REF.
  %  This is also the same as giving an \(\mathcal{O}_S\)-algebra map
  %  \(\varphi : \mathcal{A} \to f_*\mathcal{O}_T\), see
  %  Sheaves, Lemma REF.
  %  We will use all three ways of thinking about \(\varphi\),
  %  without further mention. [...]
  %  In this way we have defined a functor
  %  \begin{eqnarray}
  %  \label{equation-spec}
  %  F : \Sch^{opp} & \longrightarrow & \textit{Sets} \\
  %  T & \longmapsto & F(T) = \{\text{pairs }(f, \varphi) \text{ as above}\}
  %  \nonumber
  %  \end{eqnarray}
  %  [...]
  %  The functor \(F\) is representable by a scheme."
  \textit{Source: [Stacks Project], tag 01LQ (lemma-spec + section-spec
  functor description); [Hartshorne], II~Ex.~5.17(b).}
  Let \(X\) be a scheme and \(\mathcal{A}\) a quasi-coherent \(\mathcal{O}_X\)-algebra.
  Then \(\underline{\Spec}_X(\mathcal{A})\) represents the functor
  \(F : \Sch^{op} \to \mathrm{Set}\),
  $T \mapsto F(T) := \{(f, \varphi) : f : T \to X \text{ a morphism},\;
  \varphi : f^* \mathcal{A} \to \mathcal{O}_T \text{ an } \mathcal{O}_T\text{-algebra map}\}.$
  Equivalently, for any \(X\)-scheme \(g : T \to X\),
  there is a natural bijection
  \[
    \Hom_X(T, \underline{\Spec}_X(\mathcal{A})) \;\;\cong\;\;
    \Hom_{\mathcal{O}_X\text{-alg}}(\mathcal{A}, g_* \mathcal{O}_T)
  \]
  (using the adjunction between \(g_*\) and \(g^*\) to convert between
  \(\mathcal{A} \to g_*\mathcal{O}_T\) and \(g^*\mathcal{A} \to \mathcal{O}_T\)).
  The bijection is natural in \(T\).
\end{theorem}

% SOURCE QUOTE PROOF: [Stacks Project], tag 01LQ, proof of lemma-spec
% (read from references/stacks-constructions.tex, L553-L599).
% "We are going to use Schemes, Lemma \ref{schemes-lemma-glue-functors}.
%  First we check that $F$ satisfies the sheaf property for the
%  Zariski topology. Namely, suppose that $T$ is a scheme,
%  that $T = \bigcup_{i \in I} U_i$ is an open covering,
%  and that $(f_i, \varphi_i) \in F(U_i)$ such that
%  $(f_i, \varphi_i)|_{U_i \cap U_j} = (f_j, \varphi_j)|_{U_i \cap U_j}$.
%  This implies that the morphisms $f_i : U_i \to S$
%  glue to a morphism of schemes $f : T \to S$ such that
%  $f|_{U_i} = f_i$, see Schemes, Section \ref{schemes-section-glueing-schemes}.
%  Thus $f_i^*\mathcal{A} = f^*\mathcal{A}|_{U_i}$ and by assumption the
%  morphisms $\varphi_i$ agree on $U_i \cap U_j$. Hence by Sheaves,
%  Section \ref{sheaves-section-glueing-sheaves} these glue to a
%  morphism of $\mathcal{O}_T$-algebras $f^*\mathcal{A} \to \mathcal{O}_T$.
%  This proves that $F$ satisfies the sheaf condition with respect to
%  the Zariski topology.
%  Let $S = \bigcup_{i \in I} U_i$ be an affine open covering.
%  Let $F_i \subset F$ be the subfunctor consisting of
%  those pairs $(f : T \to S, \varphi)$ such that
%  $f(T) \subset U_i$.
%  We have to show each $F_i$ is representable.
%  This is the case because $F_i$ is identified with
%  the functor associated to $U_i$ equipped with
%  the quasi-coherent $\mathcal{O}_{U_i}$-algebra $\mathcal{A}|_{U_i}$,
%  by Lemma \ref{lemma-spec-base-change}.
%  Thus the result follows from Lemma \ref{lemma-spec-affine}.
%  Next we show that $F_i \subset F$ is representable by open immersions.
%  Let $(f : T \to S, \varphi) \in F(T)$. Consider $V_i = f^{-1}(U_i)$.
%  It follows from the definition of $F_i$ that given $a : T' \to T$
%  we have $a^*(f, \varphi) \in F_i(T')$ if and only if $a(T') \subset V_i$.
%  Finally, we have to show that the collection $(F_i)_{i \in I}$
%  covers $F$. Let $(f : T \to S, \varphi) \in F(T)$.
%  Consider $V_i = f^{-1}(U_i)$. Since $S = \bigcup_{i \in I} U_i$
%  is an open covering of $S$ we see that $T = \bigcup_{i \in I} V_i$
%  is an open covering of $T$. Moreover $(f, \varphi)|_{V_i} \in F_i(V_i)$.
%  This finishes the proof of the lemma."
\begin{proof}
  
  \leanok
  % NOTE (iter-179 review): supersedes the iter-176 placeholder-discharge note.
  % The iter-179 Block A refactor replaces the iter-176 placeholder bodies with
  % the Mathlib values ``(AffineZariskiSite.relativeGluingData _).glued`` /
  % ``.toBase``; iter-179 Lane B closes the Lean ``UniversalProperty`` proof
  % via ``isAffineHom_of_forall_exists_isAffineOpen`` +
  % ``relativeGluingData.toBase_preimage_eq_opensRange_ι`` +
  % ``isAffineOpen_opensRange`` — a real proof against the gluing infrastructure,
  % NOT against a trivial identity placeholder. The Lean type remains the
  % strictly weaker ``IsAffineHom (structureMorphism 𝒜)'' rather than the
  % Yoneda-bijection form pinned above (iter-174+ commitment; still pending
  % iter-180+). Consumers may now rely on the affineness conclusion, but the
  % Yoneda-bijection witness is still deferred.
  Following Stacks tag 01LQ (proof of lemma-spec): the functor \(F\) is a Zariski sheaf
  because morphisms of schemes and \(\mathcal{O}\)-algebra maps both glue under open covers.
  Choose an affine open cover \(X = \bigcup U_i\). The subfunctor
  \(F_i \subseteq F\) of pairs \((f, \varphi)\) with \(f(T) \subseteq U_i\) is representable by
  \(\Spec(\mathcal{A}(U_i))\) via the affine-base case
  (\cref{thm:relative_spec_affine_base}) together with the
  restriction identification \(F_i \cong F_{U_i,\,\mathcal{A}|_{U_i}}\).
  Each \(F_i \hookrightarrow F\) is an open
  subfunctor (pull back \(U_i\) along \(f\)). The \(F_i\) jointly cover \(F\) since the \(U_i\)
  cover \(X\). The gluing of representable open subfunctors produces a representing scheme,
  which is canonically isomorphic to the glued scheme \(\underline{\Spec}_X(\mathcal{A})\)
  of \cref{thm:relative_spec_exists}.
\end{proof}

\section{Affine base}
\label{sec:relspec_affine}

\begin{theorem}\leanok
[Relative spectrum over an affine base]
  \label{thm:relative_spec_affine_base}
  \lean{AlgebraicGeometry.Scheme.RelativeSpec.affine_base_iff}
  \uses{def:qc_sheaf_of_algebras}
  % NOTE (iter-173 review): the iter-173 file-skeleton encodes this as the
  % weaker ``IsAffine ((Spec R).RelativeSpec 𝒜)'' rather than the canonical
  % iso ``Spec_X(𝒜) ≅ Spec(Γ(X, 𝒜))'' pinned by this prose. The Lean name
  % ``affine_base_iff'' is also misleading (the encoded type is not an iff).
  % Iter-174+ refines to the iff-with-Γ statement once Γ : QcohAlgebra → CommRingCat
  % is available.
  
  % SOURCE: [Stacks Project], tag 01LO (lemma-spec-affine)
  % (read from references/stacks-constructions.tex, L491-L545).
  % SOURCE QUOTE:
  % "In Situation REF.
  %  Let \(F\) be the functor associated to \((S, \mathcal{A})\) above.
  %  If \(S\) is affine, then \(F\) is representable by the
  %  affine scheme \(\Spec(\Gamma(S, \mathcal{A}))\)."
  \textit{Source: [Stacks Project], tag 01LO (lemma-spec-affine).}
  Let \(X = \Spec R\) be an affine scheme and let \(\mathcal{A}\) be a quasi-coherent
  \(\mathcal{O}_X\)-algebra. Set \(A := \Gamma(X, \mathcal{A})\), regarded as an
  \(R\)-algebra. Then there is a canonical isomorphism of \(X\)-schemes
  \[
    \underline{\Spec}_X(\mathcal{A}) \;\;\cong\;\; \Spec(A).
  \]
  In particular, on an affine base the relative spectrum reduces to the absolute
  spectrum of the global sections.
\end{theorem}

% SOURCE QUOTE PROOF: "Write \(S = \Spec(R)\) and \(A = \Gamma(S, \mathcal{A})\).
% Then \(A\) is an \(R\)-algebra and \(\mathcal{A} = \widetilde A\).
% The ring map \(R \to A\) gives rise to a canonical map
% \(\)
% f_{univ} : \Spec(A)
% \longrightarrow
% S = \Spec(R).
% \(\)
% We have
% \(f_{univ}^*\mathcal{A} =  \widetilde{A \otimes_R A}\)
% by Schemes, Lemma REF.
% Hence there is a canonical map
% \(\)
% \varphi_{univ} :
% f_{univ}^*\mathcal{A} = \widetilde{A \otimes_R A}
% \longrightarrow
% \widetilde A = \mathcal{O}_{\Spec(A)}
% \(\)
% coming from the \(A\)-module map \(A \otimes_R A \to A\),
% \(a \otimes a' \mapsto aa'\). We claim that the pair
% \((f_{univ}, \varphi_{univ})\) represents \(F\) in this case."
\begin{proof}
  
  \leanok
  % NOTE (iter-179 review): supersedes the iter-176 placeholder-discharge note.
  % The iter-176 ``change IsAffine (Spec R); inferInstance`` form is now stale
  % — iter-179 Lane B closes via ``UniversalProperty 𝒜`` + ``isAffine_of_isAffineHom``
  % against the iter-179 Mathlib-backed ``RelativeSpec``/``structureMorphism``
  % bodies. The Lean type remains the strictly weaker ``IsAffine`` rather than
  % the canonical iso ``Spec_X(𝒜) ≅ Spec(Γ(X,𝒜))'' named below (iter-174+
  % commitment; still pending iter-180+).
  Write \(X = \Spec R\) and \(A = \Gamma(X, \mathcal{A})\). Quasi-coherence of
  \(\mathcal{A}\) gives \(\mathcal{A} = \widetilde{A}\) as \(\mathcal{O}_X\)-modules
  (with the \(R\)-algebra structure on \(A\) inducing the algebra structure on
  \(\widetilde A\)). The ring map \(R \to A\) gives a morphism
  \(f_{\mathrm{univ}} : \Spec A \to \Spec R = X\) and a canonical
  \(\mathcal{O}_{\Spec A}\)-algebra map
  $\varphi_{\mathrm{univ}} : f_{\mathrm{univ}}^* \mathcal{A} = \widetilde{A \otimes_R A}
  \to \mathcal{O}_{\Spec A} = \widetilde A$ given by multiplication
  \(a \otimes a' \mapsto aa'\). We verify \((\Spec A, f_{\mathrm{univ}}, \varphi_{\mathrm{univ}})\)
  represents \(F\) on the affine base: given a test scheme \(T\) and a pair
  \((f, \varphi)\), the composite
  $A = \Gamma(X, \mathcal{A}) \xrightarrow{f^\sharp}
  \Gamma(T, f^*\mathcal{A}) \xrightarrow{\varphi} \Gamma(T, \mathcal{O}_T)$
  is a ring map \(A \to \Gamma(T, \mathcal{O}_T)\) and hence corresponds to a morphism
  \(T \to \Spec A\) by the affine universal property
  (Mathlib \texttt{Scheme.toSpec\_naturality} /
  \texttt{Scheme.Hom.toSpec}). One checks this assignment inverts the universal map
  \((f, \varphi) \mapsto a\) defined by \(f = f_{\mathrm{univ}} \circ a\),
  \(\varphi = a^* \varphi_{\mathrm{univ}}\).
\end{proof}

\section{Lean encoding}
\label{sec:relspec_lean_encoding}

The Lean target file is the new \texttt{AlgebraicJacobian/Picard/RelativeSpec.lean}.
The construction proceeds by gluing absolute affine pieces: for each affine open
\(U \subseteq X\) with \(U = \Spec R\) and \(A := \mathcal{A}(U)\) regarded as an
\(R\)-algebra, the local piece is \(\Spec A\) with structure morphism \(\Spec A \to \Spec R = U\)
induced by the algebra unit \(R \to A\). Mathlib provides \texttt{IsAffineOpen.toScheme} for
the affine-open inclusion and the absolute \texttt{Scheme.Spec} functor for the local
pieces; the transition isomorphisms come from quasi-coherence of \(\mathcal{A}\) via
\texttt{QuasiCoherent.tensorIso} (or the explicit
\(\mathcal{A}(U) \otimes_R S \xrightarrow{\sim} \mathcal{A}(V)\) for \(V = \Spec S \subseteq U\)).
The gluing backbone is \texttt{AlgebraicGeometry.Scheme.GlueData} (general construction) or
the convenience wrapper \texttt{AlgebraicGeometry.AffineScheme.glueOpens}; the cocycle
condition reduces to Stacks \texttt{lemma-transitive-spec}. The universal property
(\cref{thm:relative_spec_univ}) is expressed as a
\texttt{CategoryTheory.Functor.RepresentableBy} witness, providing the natural bijection
$\Hom_X(T, \underline{\Spec}_X(\mathcal{A})) \cong
\Hom_{\mathcal{O}_X\text{-alg}}(\mathcal{A}, g_*\mathcal{O}_T)$ as a
\texttt{Yoneda}-style structure. No new core axiom is required; the entire chapter is a
definitional/structural construction on top of existing Mathlib affine-scheme
infrastructure.

\section{Out of scope}
\label{sec:relspec_out_of_scope}

This chapter packages \emph{only} the relative-spectrum construction and its core
properties. The following downstream constructions live in separate sibling chapters and
are explicitly out of scope here:
\begin{itemize}
  \item \textbf{Relative Picard functor \(\Pic^\sharp_{C/k}\)} (A.1.c) --- consumed-by, not
    defined-here; its blueprint home is the planned chapter
    \texttt{Picard\_RelPicFunctor.tex}.
  \item \textbf{Line-bundle pullback on a product \(C \times_k T\)} (A.1.b) ---
    \texttt{Picard\_LineBundlePullback.tex} (planned).
  \item \textbf{Connectedness, dimensionality, smoothness of \(\underline{\Spec}_X(\mathcal{A})\)}
    --- downstream properties that depend on assumptions on \(\mathcal{A}\)
    (locally finite type, flat, etc.), not unconditional consequences of the construction.
  \item \textbf{Identity-component \(\Pic^0\) and abelian-variety structure} (A.3) --- a
    separate Route~A sub-row, fed by but not constructed in this chapter.
\end{itemize}
